\section{Repaso de los fundamentos de los MDP}

Esta sesión la imparte el asistente de docencia Anakiette como repaso general del curso: parte de los fundamentos de los MDP y recorre la cadena teórica completa de Q-learning.

\textit{``We're going to start off by just doing a very very brief overview of MDPs.''} Esta frase inicial es a la vez una indicación y una guía para el ritmo de estudio: recuerda que, dentro de la cadena teórica completa, primero se establece el sistema de notación formal y después se profundiza en los detalles nivel por nivel.

\subsection{Procesos de decisión de Markov}

\begin{importantbox}{La quíntupla de un MDP}
Un MDP se define mediante $(\mathcal{S}, \mathcal{A}, P, R, \gamma)$:
\begin{itemize}
    \item $\mathcal{S}$: espacio de estados
    \item $\mathcal{A}$: espacio de acciones
    \item $P(s' | s, a)$: probabilidad de transición de estados
    \item $R(s, a)$ o $R(s, a, s')$: función de recompensa
    \item $\gamma \in [0, 1)$: factor de descuento
\end{itemize}
\textbf{Propiedad de Markov}: el siguiente estado depende solo del estado y la acción actuales, no del historial.
\end{importantbox}

\subsection{Función de valor y función Q}

\begin{importantbox}{Definiciones fundamentales}
\textbf{Función de valor de estado}: retorno acumulado esperado al seguir la política $\pi$ desde el estado $s$
$$V^\pi(s) = \mathbb{E}_\pi\left[\sum_{t=0}^{\infty} \gamma^t R(s_t, a_t) \mid s_0 = s\right]$$

\textbf{Función de valor de acción} (función Q): retorno acumulado esperado al ejecutar la acción $a$ en el estado $s$ y seguir después la política $\pi$
$$Q^\pi(s, a) = \mathbb{E}_\pi\left[\sum_{t=0}^{\infty} \gamma^t R(s_t, a_t) \mid s_0 = s, a_0 = a\right]$$

Relación entre ambas: $V^\pi(s) = \sum_a \pi(a|s) Q^\pi(s, a)$
\end{importantbox}

\subsection{Ecuaciones de Bellman}

\begin{importantbox}{Ecuación de esperanza de Bellman}
$$Q^\pi(s, a) = R(s, a) + \gamma \sum_{s'} P(s'|s,a) \sum_{a'} \pi(a'|s') Q^\pi(s', a')$$

\textbf{Ecuación de optimalidad de Bellman}:
$$Q^*(s, a) = R(s, a) + \gamma \sum_{s'} P(s'|s,a) \max_{a'} Q^*(s', a')$$
\end{importantbox}

\subsection{Iteración de políticas e iteración de valores}

\begin{knowledgebox}{Dos enfoques iterativos complementarios}
La iteración de políticas (policy iteration) y la iteración de valores (value iteration) pueden verse como dos formas de llevar a la práctica la ecuación de optimalidad de Bellman: la primera optimiza de manera alternada la política y la función de valor; la segunda itera directamente sobre el valor óptimo. En cuanto a la convergencia, ambas dependen de que la iteración de Bellman sea una contracción; sin embargo, la primera actualiza la política con mayor precisión en cada ronda, mientras que la segunda requiere un cálculo más sencillo en cada paso.
\end{knowledgebox}

\begin{table}[H]
\centering
\begin{tabular}{ll}
\toprule
Método & Características \\
\midrule
Iteración de políticas & Construye explícitamente la política de comportamiento (alineada con la política anterior) y cambia la política solo cuando la función de valor converge \\
Iteración de valores & Itera directamente $V(s) \leftarrow \max_a [R+\gamma V(s')]$, sin almacenar una política explícita \\
Relación entre ambas & La iteración de valores puede verse como un caso particular de la iteración de políticas (la política se extrae solo al final) \\
\bottomrule
\end{tabular}
\caption{Comparación de los operadores iterativos más comunes en MDP discretos}
\label{tab:mdp-iteration}
\end{table}

\subsection{Resumen de la sección}

Los MDP proporcionan el marco formal del RL. La función de valor y la función Q son las herramientas fundamentales para evaluar la calidad de una política, y las ecuaciones de Bellman describen la relación recursiva entre ellas.
